\documentclass[12pt]{article}

\usepackage[utf8]{inputenc}
\usepackage[T1]{fontenc}
\usepackage[romanian]{babel}
\usepackage{amsmath, amssymb, amsthm}
\usepackage{geometry}
\geometry{a4paper, margin=2.5cm}

\title{\Large\textbf{Teorema de aproximare a lui Runge}}
\date{}

\theoremstyle{plain}
\newtheorem{theorem}{Teoremă}

\begin{document}

\maketitle

\begin{theorem}[Runge]
Fie $\Omega \subset \mathbb{C}$ o mulțime deschisă și fie $K \subset \Omega$ o submulțime compactă
astfel încât $\mathbb{C} \setminus K$ este conexă.

Atunci pentru orice funcție $f$ olomorfă în $\Omega$ și pentru orice $\varepsilon > 0$ există un polinom $P$
cu proprietatea


\[
\sup_{z \in K} |f(z) - P(z)| < \varepsilon.
\]



Cu alte cuvinte, orice funcție olomorfă pe o mulțime deschisă poate fi aproximată uniform pe compacte
prin polinoame, atâta timp cât complementul compactului este conex.
\end{theorem}

\end{document}
